% A matrix product state: arrays A on the chain, bonds between them, occupation
% indices n_i down -- a finite basis, so no wavy legs.
%
% The same chain twice, in the two conventions a tensor is drawn in. Square is
% this notation's default, because elsewhere it tells an array of numbers from a
% function of position, which is a circle. Here there are no functions, so there
% is nothing to tell apart, and a great deal of the literature draws every
% tensor of a network as a circle -- that is the lower row, which is
% `{coef, circle}' and nothing else: same style, same colour, same meaning.
%
% Shape does carry meaning once a chain is in canonical form, where a triangle
% says isometric and a diamond says centre (04, 05). That is a different claim
% from this one, and it is not what `circle' touches.
%
% The fourth slot is the dots of a chain that goes on: the bonds either side of
% it stop at its border.
\documentclass[border=4pt]{standalone}
\usepackage{amsmath,amssymb}
\usepackage{tikz-tensors}
\input{notation}
\begin{document}
\begin{tikzpicture}
  \tnstack{P}{5}{ket}
  \tnlayer{P}{ket}{coef/$A$, coef/$A$, coef/$A$, dots, coef/$A$}
  \tnconnect{P}
  \tnopen{down}{P}
  \tnmid{P-ket-1}{P-ket-2}{$D$}
  \foreach \i/\n in {1/1, 2/2, 3/3, 5/N} \tnput[below]{P-\i-down}{$n_{\n}$};
  \tnbreak
  \tnstack{Q}{5}{ket}
  \tnlayer{Q}{ket}{{coef, circle}/$A$, {coef, circle}/$A$, {coef, circle}/$A$,
                   dots, {coef, circle}/$A$}
  \tnconnect{Q}
  \tnopen{down}{Q}
  \tnmid{Q-ket-1}{Q-ket-2}{$D$}
  \foreach \i/\n in {1/1, 2/2, 3/3, 5/N} \tnput[below]{Q-\i-down}{$n_{\n}$};
\end{tikzpicture}
\end{document}
